\begin{abstract}
Large language models (LLMs) are trained on human language corpora whose hierarchical structure induces pronounced anisotropy in learned representations. Across layers and training stages, activations and parameters exhibit a spike-tail spectrum: a small number of directions concentrate disproportionate energy, while the remaining directions form a broad long tail.

This anisotropy constitutes a structural curse for large-scale training. High-energy components inflate elementwise dynamic range, generate extreme activation outliers, and distort adaptive second-moment statistics. In low-bit regimes, dominant spike directions control quantization scale and compress long-tail semantic variation. Mitigating these effects typically requires explicit spectral methods such as SVD, which are computationally expensive and hardware-unfriendly.

We show that this anisotropy is not accidental but emerges from a fundamental inductive bias of language modeling: the meaning of a token is determined by its context. Self-attention performs contextual aggregation through convex combinations, and when coupled with non-odd nonlinearities and residual accumulation, systematically reinforces shared semantic components. In high dimensions, these shared components accumulate into coherent mean directions. We demonstrate that this mean bias dominates both extreme elementwise activations and optimizer second-moment inflation.

Crucially, this structure is also a blessing. Because the mean bias-induced dominant instability is rank-one, it can be isolated with a simple mean-subtraction operation. This bias-centric conditioning achieves the core benefits of spectral control without costly decompositions. We provide theoretical analysis explaining mean emergence and dominance, and empirically validate the approach in large-scale LLM training. Across low-bit quantization and adaptive optimization settings, source-level mean removal substantially reduces instability while preserving semantic covariance.
\end{abstract}

\section{Introduction}

Large language models are trained on massive human language corpora, which induce a characteristic geometry in learned representations. Across layers and training stages, activations, gradients, and model parameters exhibit pronounced anisotropy: a small number of directions concentrate disproportionate energy (strong spectral spikes), while the remaining directions form a broad long tail \citep{ethayarajh2019, mu2018, timkey2021, kovaleva2021, machina2024}.

While linguistically natural, this structure becomes a computational liability during large-scale training. Anisotropy inflates elementwise dynamic range, amplifies extreme activation outliers, and distorts adaptive second-moment statistics \citep{kingma2014, loshchilov2017, huang2026}. In low-bit regimes, dominant spike directions control quantization scale, compressing long-tail variation into narrow numerical bins \citep{xiao2023, frantar2023, lin2024, li2024, ashkboos2024}. To mitigate these effects, prior approaches rely on explicit spectral control, often via singular value decomposition (SVD) or related orthogonalization procedures, as in Metis \citep{cao2025}, Muon \citep{jordan2024}, and Spectra \citep{huang2026}. Although effective, such techniques are computationally expensive, hardware-unfriendly, and difficult to integrate into large-scale training loops. In this sense, anisotropy constitutes a structural curse for efficient model training.

In this work, we examine the internal structure of anisotropy and uncover a sharp asymmetry: anisotropy in activations is dominated by a coherent mean bias. This phenomenon is not accidental. At every representational level---tokens, phrases, clauses, and full sentences---language exhibits shared contextual signal within a given context. Contextual aggregation mechanisms such as self-attention, together with nonlinear transformations, integrate these shared components across tokens. This aggregation naturally induces coherent mean directions in the representation space. As deeper layers capture progressively higher-level semantic regularities, such coherent components reappear and accumulate, producing layer-wise mean drift that propagates through residual connections.

Importantly, due to high dimensionality, even moderate coherent bias produces a consistent directional drift across thousands of coordinates. This rank-one component dominates elementwise extreme values and inflates adaptive second-moment statistics, despite accounting for a small fraction of total variance, making it the primary driver of outlier formation and optimizer instability. Crucially, this structure is not merely harmful---it is also exploitable. Because the dominant instability arises from a rank-one mean component, it can be isolated and removed with a single mean-subtraction operation, which is hardware-friendly and dramatically cheaper than repeated layer-wise spectral decompositions. The same structure that produces instability also exposes a low-dimensional handle for mitigation: the blessing of bias.

We validate this bias-centric interpretation in two downstream settings. In low-bit quantization, mean removal substantially reduces dynamic-range inflation and improves long-tail resolution \citep{xiao2023, wei2023, ashkboos2024, cao2025}. In adaptive optimization, it contracts extreme second-moment tails, leading to more balanced effective learning-rate allocation \citep{kingma2014, huang2026}. These results indicate that addressing anisotropy at its structural origin can replace or complement heavy spectral interventions with minimal computational overhead.

\paragraph{Contributions.}
\begin{enumerate}
    \item We identify coherent activation mean bias as a primary structural source of anisotropy in LLM training.
    \item We establish a forward/backward asymmetry: activations accumulate coherent mean drift, while gradient outputs remain mean-neutral.
    \item We quantify that mean-dominated components account for the majority of extreme activation magnitudes and second-moment inflation.
    \item We propose source-level mean removal with optional spike whitening as efficient spectral conditioning mechanisms.
    \item We demonstrate consistent improvements in low-bit quantization and adaptive optimization.
\end{enumerate}

\section{Structural Emergence of Mean Bias in Transformer Representations}
\label{sec:structural_emergence}

This section demonstrates that the coherent mean bias observed in Transformer models is not an accidental artifact but a predictable consequence of the interaction between language statistics and the specific geometry of high-dimensional operators. We argue that this bias emerges through a three-stage causal chain: frequency-weighted initialization, architectural regeneration, and residual accumulation.

\paragraph{Statistical Origin: Frequency-Weighted Embedding Expectation.}
A ``seed'' bias is introduced at the very first layer of the model. Let $\mathcal{V}$ denote the vocabulary and let $p(v)$ be the token frequency distribution over $\mathcal{V}$. In natural language corpora, $p(v)$ is highly non-uniform and typically follows a Zipfian distribution. The expected embedding vector over the corpus is:
\begin{equation}
\mu_{\text{embed}} = \mathbb{E}_{v \sim p}[E_v] = \sum_{v \in \mathcal{V}} p(v) E_v,
\end{equation}
where $E_v \in \mathbb{R}^H$ denotes the embedding vector of token $v$ \citep{mikolov2013a, arora2017}.
Because high-frequency tokens (e.g., articles, punctuation, and syntactic markers) dominate $p(v)$ and receive disproportionately many gradient updates during training, their embedding vectors tend to align along common predictive directions. In high-dimensional spaces, such coherent alignment induces a dominant rank-one component in the embedding matrix even before subsequent Transformer processing begins.

\paragraph{The Mechanism of Mean Regeneration.}
Even if internal activations are perfectly centered at zero, the fundamental operators of the Transformer architecture systematically regenerate a coherent mean. This regeneration is not incidental but arises from the algebraic structure of nonlinearities, attention mechanisms, and residual integration.

\textbf{Non-Odd Nonlinearities as Bias Generators:} Most modern activation functions $\phi(\cdot)$ used in Transformer architectures---such as ReLU, GELU, and SwiGLU---are non-odd functions. For any symmetric zero-mean random variable $z$ satisfying $\mathbb{E}[z]=0$ and $\mathbb{P}(z \neq 0) > 0$, non-odd nonlinearities generally satisfy:
\begin{equation}
\mathbb{E}[\phi(z)] > 0,
\end{equation}
under mild regularity conditions \citep{timkey2021, puccetti2022}.
For example, ReLU is defined as $\phi(z) = \max(0,z)$. If $z$ is symmetric, then:
\begin{equation}
\mathbb{E}[\max(0,z)] = \int_0^\infty z \, p(z) dz > 0,
\end{equation}
because all negative mass is truncated. Thus, ReLU introduces a strictly positive drift even when the input is perfectly centered. The same phenomenon holds approximately for GELU and SwiGLU, which suppress negative regions asymmetrically.\footnote{While SwiGLU has a small negative basin, its expectation over a symmetric input remains strictly positive due to its dominant positive linear asymptote.} Only strictly odd nonlinearities such as $\tanh$ preserve a zero mean in expectation. Therefore, unless $\phi(\cdot)$ is explicitly odd, mean regeneration is mathematically unavoidable.

\textbf{Softmax Attention as Convex Aggregation:} Self-attention further reinforces mean drift. Given query, key, and value projections $Q = XW_Q$, $K = XW_K$, $V = XW_V$, the attention output is $Y_i = \sum_j A_{ij} V_j$, where $A = \mathrm{softmax}(QK^\top / \sqrt{d})$. Because softmax enforces positivity and row-wise normalization ($A_{ij} \ge 0$ and $\sum_j A_{ij} = 1$), each output $Y_i$ is a convex combination of value vectors. Decompose the values as:
\begin{equation}
V_j = c + \epsilon_j,
\end{equation}
where $c$ is a shared contextual component and $\epsilon_j$ are zero-mean deviations ($\sum_j \epsilon_j \approx 0$). Then:
\begin{equation}
Y_i = \sum_j A_{ij} (c + \epsilon_j) = c + \sum_j A_{ij} \epsilon_j.
\end{equation}
The shared component $c$ is preserved exactly. Meanwhile, the deviations $\epsilon_j$ are averaged. Because $A_{ij}$ depends on $QK^\top$ (which itself depends on $V$ through previous layers), the strict symmetry conditions required for the deviations to cancel perfectly are violated. Thus, attention acts as a data-dependent low-pass filter, preserving coherent components while attenuating fluctuations.

\textbf{Residual Integration and High-Dimensional Magnification:} Mean drift generated by nonlinearities and attention is integrated through residual connections:
\begin{equation}
\mu_{l+1} = \mu_l + \Delta \mu_l,
\end{equation}
where $\Delta \mu_l$ denotes the mean generated by layer $l$. This additive coupling causes bias to accumulate approximately linearly with depth when $\Delta \mu_l$ maintains consistent orientation.

Crucially, high dimensionality magnifies this effect. Let $\bar{\mu}$ denote the average coordinate-wise bias. The Euclidean norm of the mean vector scales as:
\begin{equation}
\|\mu\|_2 \sim \sqrt{H} \bar{\mu},
\end{equation}
where $H$ is the hidden dimension \citep{tropp2012, halko2011}. In a space with $H=4096$, a seemingly small per-coordinate bias is amplified by a factor of $\sqrt{4096} = 64$ in total magnitude. This geometric magnification explains why mean-dominated components account for a large fraction (empirically 85\%--90\%) of extreme activation outlier energy.

\paragraph{The LayerNorm Caveat: Axes of Expectation.}
A rigorous reader might ask: does layer normalization (LayerNorm or RMSNorm) not eliminate this mean drift? Crucially, LayerNorm centers activations strictly across the \textit{hidden dimension} for a single token ($\sum_{h=1}^H x_{ih} = 0$). However, the coherent drift $\mu$ we analyze here occurs across the \textit{batch or sequence distribution} ($\mathbb{E}_i[x_{ih}] = \mu_h \neq 0$). Because LayerNorm does not center across tokens, the token population can maintain a massive expected value in a specific coordinate, so long as the sum of coordinates for each token is zero ($\sum_{h=1}^H \mu_h = 0$). This allows the structural bias to persist and amplify.

\paragraph{Empirical Verification.} We measure mean magnitude using the normalized ratio $R = \frac{\|\mu\|_2}{\sqrt{\mathbb{E}\|X\|_F^2/N}}$. Across layers and training stages, we observe:
\begin{itemize}
    \item Activation $R$ ranges from 10 to 200+.
    \item Gradient-output $R \approx 1$ (mean-neutral).
    \item Mean bias strengthens over training and increasingly dominates the top 0.1\% of activation magnitudes.
\end{itemize}

\section{Mean Bias as the Dominant Source of Activation Outliers}
\label{sec:activation_outliers}

While Section~\ref{sec:structural_emergence} detailed the architectural regeneration of bias, we now address the sharp structural consequence: the production of extreme elementwise activations that destabilize low-bit quantization and adaptive optimization. Because the dynamic range of low-bit quantization is governed by entrywise maxima, we perform an elementwise mechanical decomposition to locate the source of instability.

To attribute outlier energy, we decompose the activation matrix $X \in \mathbb{R}^{N \times H}$ into three mutually orthogonal components:
\begin{equation}
X = \underbrace{\mathbf{1}\mu^\top}_{M (\text{mean})} + \underbrace{U_k \Sigma_k V_k^\top}_{X_{\text{spike}}} + \underbrace{X_{\text{tail}}}_{\text{residual}},
\end{equation}
where $\mu = \frac{1}{N}X^\top \mathbf{1}$, $X_{\text{spike}}$ is the projection of the centered matrix onto the top $k = 1\% \cdot H$ singular directions, and $X_{\text{tail}}$ is the remaining orthogonal residual. The decomposition is orthogonal in the Frobenius inner product, hence $\|X\|_F^2 = \|M\|_F^2 + \|X_{\text{spike}}\|_F^2 + \|X_{\text{tail}}\|_F^2$.

\paragraph{Outlier Attribution.} 
Define the outlier set $\mathcal{E}_{\text{top}}$ as the top 0.1\% entries ranked by magnitude $|X_{ij}|$. For each $(i,j) \in \mathcal{E}_{\text{top}}$, we define squared-magnitude shares:
\begin{equation}
\rho_{ij}^{(\text{mean})} = \frac{M_{ij}^2}{X_{ij}^2}, \quad \rho_{ij}^{(\text{spike})} = \frac{(X_{\text{spike}})_{ij}^2}{X_{ij}^2}, \quad \rho_{ij}^{(\text{tail})} = \frac{(X_{\text{tail}})_{ij}^2}{X_{ij}^2}.
\end{equation}

Table 1 summarizes statistics across layers and training stages. Two structural transitions are visible. In shallow layers early in training, mean bias already accounts for the majority of extreme entries. As training proceeds, shallow-layer outliers become almost entirely mean-dominated. In deeper layers, variance initially dominates outliers, but by late training, mean bias becomes the principal driver there as well.

\paragraph{High-Dimensional Extreme-Value Amplification.} 
We now mathematically explain why a rank-one mean inevitably dominates extremes in high dimension.

\begin{theorem}[High-Dimensional Extreme Dominance]
Let $X_{ij} = \mu_j + Z_{ij}$, where $\mu_j$ is a deterministic mean shift and $Z_{ij}$ is independent zero-mean sub-Gaussian noise with variance proxy $\sigma^2$. For a high quantization threshold $t < |\mu_j|$, the probability of the coordinate exceeding the threshold is bounded from below by:
\begin{equation}
\mathbb{P}(|X_{ij}| > t) \ge \mathbb{P}(|Z_{ij}| < |\mu_j| - t).
\end{equation}
\end{theorem}

\begin{proof}
By the reverse triangle inequality, $|X_{ij}| = |\mu_j + Z_{ij}| \ge |\mu_j| - |Z_{ij}|$. Therefore, if the noise satisfies $|Z_{ij}| < |\mu_j| - t$, it strictly follows that $|X_{ij}| > t$. Because $Z_{ij}$ is concentrated around zero, the probability $\mathbb{P}(|Z_{ij}| < |\mu_j| - t)$ approaches $1$ as $\sigma$ becomes small relative to the mean shift.

In stark contrast, if a coordinate is purely variance-driven ($X'_{ij} = Z_{ij}$ without a mean shift), exceeding the threshold $t$ requires a rare tail event, which is bounded by standard sub-Gaussian concentration:
\begin{equation}
\mathbb{P}(|Z_{ij}| > t) \le 2\exp\left(-\frac{t^2}{2\sigma^2}\right).
\end{equation}
\end{proof}

This result formalizes a key statistical principle: a deterministic mean shift generates extreme values with constant probability $\mathcal{O}(1)$, while purely variance-driven extremes occur with exponentially vanishing probability $\mathcal{O}(\exp(-t^2/\sigma^2))$. Since quantization scale, clipping thresholds, and adaptive second-moment updates are governed by scalar extreme values, mean bias becomes the first-order structural driver of dynamic-range inflation, second-moment distortion, and long-tail semantic compression. Mitigating mean bias therefore directly targets the principal geometric source of extreme-value instability.


\section{Mean Bias as the Dominant Source of Activation Outliers}
\label{sec:activation_outliers}

While Section~\ref{sec:structural_emergence} detailed the architectural regeneration of bias, we now address its sharp structural consequence: the production of extreme elementwise activations that destabilize low-bit quantization and adaptive optimization. Because the dynamic range of low-bit quantization is strictly governed by entrywise maxima (the $L_\infty$ norm), we perform an elementwise mechanical decomposition to precisely locate the geometric source of this instability.

To attribute outlier energy, we decompose the activation matrix $X \in \mathbb{R}^{N \times H}$ into three mutually orthogonal components. Let $\mu = \frac{1}{N}X^\top \mathbf{1} \in \mathbb{R}^H$ be the feature-wise mean vector, and let $\tilde{X} = X - \mathbf{1}\mu^\top$ be the column-centered activation matrix. We decompose $X$ as:
\begin{equation}
X = \underbrace{\mathbf{1}\mu^\top}_{M (\text{mean})} + \underbrace{U_k \Sigma_k V_k^\top}_{X_{\text{spike}}} + \underbrace{X_{\text{tail}}}_{\text{residual}},
\end{equation}
where $X_{\text{spike}}$ is the truncated rank-$k$ singular value decomposition of the centered matrix $\tilde{X}$ onto its top $k = \lfloor 0.01 H \rfloor$ singular directions, and $X_{\text{tail}} = \tilde{X} - X_{\text{spike}}$ is the remaining orthogonal residual. Because $X_{\text{spike}}$ and $X_{\text{tail}}$ are linear projections of the centered matrix, they maintain zero column means (i.e., they are orthogonal to the all-ones vector $\mathbf{1}$). Therefore, the decomposition is strictly orthogonal in the Frobenius inner product, yielding exact energy separation:
\begin{equation}
\|X\|_F^2 = \|M\|_F^2 + \|X_{\text{spike}}\|_F^2 + \|X_{\text{tail}}\|_F^2.
\end{equation}

\paragraph{Outlier Attribution.} 
Define the outlier set $\mathcal{E}_{\text{top}}$ as the indices $(i,j)$ of the top $0.1\%$ of entries ranked by absolute magnitude $|X_{ij}|$. To determine which geometric component drives these extreme values, we define the squared-magnitude share for each $(i,j) \in \mathcal{E}_{\text{top}}$:
\begin{equation}
\rho_{ij}^{(\text{mean})} = \frac{M_{ij}^2}{X_{ij}^2}, \quad \rho_{ij}^{(\text{spike})} = \frac{(X_{\text{spike}})_{ij}^2}{X_{ij}^2}, \quad \rho_{ij}^{(\text{tail})} = \frac{(X_{\text{tail}})_{ij}^2}{X_{ij}^2}.
\end{equation}
\textit{Remark:} While the decomposition is strictly orthogonal over the full matrix, individual elementwise squares technically contain minor cross-terms (e.g., $2M_{ij}(X_{\text{spike}})_{ij}$). Nevertheless, measuring these primary component ratios provides a highly accurate and stable proxy for identifying the dominant driver of extreme values.

Table~\ref{tab:outlier_shares} summarizes these statistics across layers and training stages. Two distinct structural transitions emerge. In shallow layers early in training, mean bias already accounts for the vast majority of extreme entries. As training proceeds, shallow-layer outliers become almost entirely mean-dominated. In deeper layers, the spike variance initially dominates outliers, but by late training, mean bias transitions into the principal driver across the entire network.

\paragraph{High-Dimensional Extreme-Value Amplification.} 
We now mathematically formalize why a rank-one mean shift inevitably dominates extreme elementwise values in high-dimensional settings, overpowering the variance.

\begin{theorem}[Elementwise Extreme Dominance]
\label{thm:extreme_dominance}
Let an activation coordinate be modeled as $X_{ij} = \mu_j + Z_{ij}$, where $\mu_j$ is a deterministic mean shift and $Z_{ij}$ is independent zero-mean sub-Gaussian noise with variance proxy $\sigma^2$. For a high quantization threshold $t > 0$ such that $t < |\mu_j|$, the probability of the coordinate magnitude exceeding the threshold satisfies:
\begin{equation}
\mathbb{P}(|X_{ij}| > t) \ge 1 - 2\exp\left(-\frac{(|\mu_j| - t)^2}{2\sigma^2}\right).
\end{equation}
\end{theorem}

\begin{proof}
By the reverse triangle inequality, $|X_{ij}| = |\mu_j + Z_{ij}| \ge |\mu_j| - |Z_{ij}|$. Therefore, if the ambient noise fluctuation satisfies $|Z_{ij}| < |\mu_j| - t$, it strictly guarantees that $|X_{ij}| > t$. 

Because $Z_{ij}$ is $\sigma^2$-sub-Gaussian, the probability of the complementary tail event is bounded by $\mathbb{P}(|Z_{ij}| \ge |\mu_j| - t) \le 2\exp\left(-\frac{(|\mu_j| - t)^2}{2\sigma^2}\right)$. Subtracting this from $1$ yields the lower bound. As the noise variance $\sigma^2$ becomes small relative to the margin $(|\mu_j| - t)^2$, this probability rapidly approaches $1$.

In stark contrast, if an activation coordinate is purely variance-driven ($X'_{ij} = Z_{ij}$ without a mean shift), exceeding the threshold $t$ requires a rare tail event, which is strictly bounded by standard sub-Gaussian concentration:
\begin{equation}
\mathbb{P}(|Z_{ij}| > t) \le 2\exp\left(-\frac{t^2}{2\sigma^2}\right).
\end{equation}
\end{proof}

This theorem formalizes a critical statistical divergence: a deterministic mean shift generates extreme values with constant probability $\mathcal{O}(1)$ (with respect to $H$), while purely variance-driven extremes occur with exponentially vanishing probability $\mathcal{O}(\exp(-t^2/\sigma^2))$. 

Because the dynamic range for blockwise low-bit tensor quantization is determined by the $L_\infty$ norm (the maximum scalar element in a block), the deterministic mean bias structurally stretches the quantization grid. This forces the grid to accommodate the rank-one drift, fatally starving the variance-driven semantic tail of bit-depth. Mitigating this mean bias therefore directly targets the geometric root of extreme-value instability, restoring resolution to the semantic signal.